\documentclass[../../../../main.tex]{subfiles}
\begin{document}

\begin{flushright}
\footnotesize\textit{【自動文字起こし・要確認】}
\end{flushright}

{\bf [解]}
$x$にある石に対し試行を行い，後の座標を$f(x)$とすると
\begin{align*}
f(x)=\begin{cases}-x&(\text{確率}\frac12)\\2-x&(\text{確率}\frac12)\end{cases}
\end{align*}
である．

\textbf{(1)} 2回試行を繰り返すと，以下のようになる．
\begin{itemize}
\item $x\to-x\to x$（確率$\frac12\cdot\frac12=\frac14$）
\item $x\to-x\to2+x$（確率$\frac14$）
\item $x\to2-x\to x-2$（確率$\frac14$）
\item $x\to2-x\to x$（確率$\frac14$）
\end{itemize}

よって，もとめる石置率$P=2\cdot\dfrac14=\dfrac12$．

\textbf{(2)} $2k$回目の試行の後（$k\in\mathbb{N}$）の石の位置を$a_k$とすると，(1)から，
\begin{align*}
a_{k+1}=\begin{cases}a_k&(\text{確率}\frac12)\\a_k+2&(\text{確率}\frac14)\\a_k-2&(\text{確率}\frac14)\end{cases}
\end{align*}
となる．$a_{k+1}=a_k-2$となる試行が1回でも行われた場合，$\max a_n=2(n-1)-2=2n-4$となるから，$a_n=2n-2$は不適である．同様に考えて，$a_n=2n-2$となるのは，$a_{k+1}=a_k$が$1$回，$a_{k+1}=a_k+2$が$n-1$回行われた時であり，もとめる確率は
\begin{align*}
P={}_n\mathrm{C}_1\left(\frac12\right)\left(\frac14\right)^{n-1}=n\left(\frac12\right)^{2n-1}．
\end{align*}

\newpage

\end{document}
